\section{Benchmark Setup}

Our goal is to test whether a compact DTI model panel preserves its ranking across synthetic cold-start regimes and real distribution shifts. The panel contains three models: \texttt{base}, a plain encoder that serves as the reference system; \texttt{RAICD}, a retrieval-augmented model that aggregates support from similar seen entities; and \texttt{FTM sparse}, a functional target-memory model designed to improve target-cold generalization. We intentionally keep the panel small so that benchmark effects are not obscured by a large architecture zoo.

We evaluate this panel on two benchmark groups. The first is a synthetic reference panel based on \texttt{BindingDB\_Kd}, with three split definitions: \texttt{unseen\_drug}, \texttt{blind\_start}, and \texttt{unseen\_target}. The second is a real-OOD screening panel. Our primary real-OOD benchmark is the BindingDB patent split (\texttt{patent\_temporal}), which uses a year-based temporal partition and is evaluated over three seeds. We also include \texttt{BindingDB\_Ki / unseen\_target} and \texttt{DAVIS / unseen\_target} as supporting screens that test whether the synthetic \texttt{unseen\_target} gain of \texttt{FTM sparse} transfers across assay and dataset shifts. Table~\ref{tab:taxonomy} summarizes the shift taxonomy and split statistics.

\begin{table}[t]
  \centering
  \scriptsize
  \caption{Benchmark taxonomy and seed0 split statistics. Test prevalence differs substantially across benchmark families, so \AUPRC{} should be compared within benchmark rather than across rows.}
  \label{tab:taxonomy}
  \setlength{\tabcolsep}{4pt}
  \resizebox{\linewidth}{!}{
  \begin{tabular}{@{}llrrr@{}}
    \toprule
    Benchmark & Shift family & Synth.? & Test rows & Pos. rate \\
    \midrule
    \texttt{Kd / unseen\_drug} & drug-cold & yes & 12,744 & 0.1966 \\
    \texttt{Kd / blind\_start} & drug+target-cold & yes & 2,821 & 0.2031 \\
    \texttt{Kd / unseen\_target} & target-cold & yes & 12,052 & 0.2239 \\
    \texttt{patent / patent\_temporal} & temporal+source & no & 49,028 & 0.5745 \\
    \texttt{Ki / unseen\_target} & assay+target-cold & mixed & 70,380 & 0.4652 \\
    \texttt{DAVIS / unseen\_target} & dataset+target-cold & mixed & 5,168 & 0.0615 \\
    \bottomrule
  \end{tabular}}
\end{table}


We report \AUPRC{} as the primary metric and \AUROC{} as a secondary metric. \AUPRC{} is emphasized because class balance and ranking behavior are central to the benchmark claims in this paper. For the synthetic reference panel and the promoted patent benchmark, we report matched three-seed results. For \texttt{BindingDB\_Ki} and \texttt{DAVIS}, we report seed0 support only and treat these as supporting rather than headline evidence. We further add paired bootstrap uncertainty support for the key ranking reversals: \texttt{RAICD} versus \texttt{base} on synthetic \texttt{blind\_start}, and \texttt{base} versus \texttt{RAICD} or \texttt{FTM} on patent temporal.

The patent benchmark requires several setup choices that are documented explicitly. We use a local year-based split of \texttt{BindingDB\_patent}, with training years before the validation year, a held-out validation year, and a later temporal test set. Because this benchmark also changes provenance and assay composition, we describe it as temporal/provenance OOD rather than pure time shift. We also use a patent-specific affinity conversion rather than reusing the standard \texttt{pKd} transform, because the patent benchmark stores affinities in a different log-scale convention. BindingDB-derived panels are anchored in the BindingDB resource \citep{liu2007bindingdb}, while DAVIS provides an external kinase-focused support benchmark \citep{davis2011kinase}.

Finally, for diagnostic analysis, we use the saved overlap score from each test prediction to compute overlap buckets that summarize how strongly each test example resembles the training support. We report benchmark taxonomy and split statistics, including test-set size and prevalence, so that \AUPRC{} is interpreted within benchmark rather than across benchmark families.
